%\renewcommand{\thefootnote}{\arabic{footnote}}
\chapter{PENDAHULUAN}
\label{BAB1:pendahuluan}

\section{Latar Belakang}
 Seiring berjalan waktu kemajuan Natural Language Processing (NLP), konsep Retrieval-Augmented Generation (RAG) mulai banyak diaplikasikan untuk menjawab permasalahan akses dan interpretasi dokumen teks yang kompleks. 
 Konsep ini pertama kali diperkenalkan oleh Lewis et al. (2023), yang menggabungkan metode retrieval berbasis pencarian semantik dengan generation menggunakan Large Language Model (LLM). Melalui gabungan tersebut hasil 
 sistem RAG dapat memberikan jawaban yang bersumber langsung dari dokumen rujukan, sehingga bisa meningkatkan akurasi dan transparansi dalam jawaban.

 Penerapan teknologi tersebut sangat relevan mengingat pengelolaan dokumen regulasi universitas merupakan hal penting dalam tata kelola akademik dan administratif. Dokumen seperti peraturan rektor, pedoman akademik maupun 
 kebijakan fakultas bersifat dinamis dan terus ada perubahan seiring dengan perkembangan kebijakan pendidikan tinggi. Perubahan tersebut bisa berupa revisi parsial, pembaruan pasal tertentu saja atau penggatian keseluruhan dokumen. 
 Namun, permasalahan sistem pencarian dokumen regulasi yang umum digunakan berbasis kata kunci (keyword-based retrieval), sehingga sering terjadi gagal membedakan versi dokumen masih berlaku dan tidak bisa menelusuri hubungan antarversi regulasi yang saling terkait. Hal ini menimbulkan kesulitan bagi civitas akademik dalam memastikan validitas dan relevansi suatu kebijakan universitas.

 Penelitian yang sudah ada menunjukan bahwa RAG memiliki kemampuan adaptif diberbagai ranah. Zhang et al. (2023) mengimplementasikan RAG pada bidang pendidikan (CHAX: A RAG-based Chatbot for Chax Education), sementara Amazou et al. (2025) menerapkannya dalam konteks hukum kontrak untuk meningkatkan akurasi dan kemampuan citation dalam sistem hukum digital.
 
 Meskipun dari hasil penelitian yang sudah memberikan hasil yang menjanjikan, tetapi penelitian tersebut masih berfokus pada dokumen statis ayng tidak memiliki hubungan versi atau konteks temporal yang bisa berubah. 
 Dalam konteks dokumen regulasi universitas, pendekatan RAG konvesional menghadapi masalah tambahan karena sifat dokumen yang amandatif (dapat berubah sebagian) dan temporal (memiliki periode berlaku). 
 Kompleksitas ini memiliki permasalahan yang perlu diselesaikan dalam pengembangan asisten digital berbasis RAG untuk regulasi universitas, yaitu:

\begin{enumerate}
    \item \textbf{Temporal Inconsistency} \\
    Sistem tidak dapat membedakan konteks waktu dari informasi, sehingga berpotensi memberikan jawaban berdasarkan dokumen yang sudah tidak berlaku.

    \newpage

    \item \textbf{Partial Supersession} \\
    Perubahan regulasi diuniversitas sifat parsial, dimana hanya beberapa pasal atau ayat yang diperbarui. Sistem RAG tanpa mekanisme pelacakan versi cenderung mencampur versi lama dan baru.

    \item \textbf{Citation Accuracy} \\
    Ketika sistem tidak memiliki manajemen versi dengan baik, refrensi terhadap sumber regulasi menjadi tidak akurat. Hal ini bisa menimbulkan resiko kesalahan interprentasi terhadap kebijakan yang berlaku.

    \item \textbf{Structural Semantics Loss} \\
    Dokumen regulasi memiliki struktur yang kompleks jika proses \textit{chunking} dan \textit{embedding} tidak mempertimbangkan struktur sistem akan kehilangan konteks semantik antarbagian dan menurukan semantic retrieval.
\end{enumerate}

\section{Rumusan Masalah}
    Berdasarkan latar belakang permasalahan yang telah diuraikan, rumusan masalah dalam penelitian ini adalah:
    \begin{enumerate}
        \item Bagaimana menerapkan metode \textit{chunking} dan \textit{embedding} yang efektif untuk mempertahankan struktur hirarki dokumen regulasi (bab, pasal, ayat) agar konteks semantik tidak hilang saat proses \textit{retrieval}?

        \item Bagaimana mekanisme \textit{retrieval} dan \textit{generation} dapat dioptimalkan untuk memastikan akurasi sitasi, sehingga jawaban yang dihasilkan sistem merujuk secara tepat pada dokumen yang relevan?

        \item Sejauh mana peningkatan kinerja sistem RAG dengan pendekatan \textit{version-aware} dibandingkan RAG konvensional dalam aspek akurasi jawaban dan relevansi dokumen regulasi universitas?
    \end{enumerate}
 
\section{Batasan Masalah}
    Agar penelitian ini lebih terarah dan fokus pada tujuan yang ingin dicapai, penulis menetapkan beberapa batasan masalah sebagai berikut:
    \begin{enumerate}
        \item \textbf{Lingkup Data}: Data yang digunakan dalam penelitian ini terbatas pada dokumen regulasi internal Universitas Pertamina (seperti Peraturan Rektor, SK, dan Pedoman Akademik) dalam format teks digital (PDF).

        \item \textbf{Model Bahasa}: Penelitian ini menggunakan \textit{Pre-trained Large Language Model} (LLM) sebagai generator jawaban (misalnya model berbasis GPT atau Llama) dan tidak melakukan proses \textit{fine-tuning} atau pelatihan ulang pada model tersebut.

        \item \textbf{Fokus Versioning}: Kemampuan \textit{version-aware} pada sistem dibatasi pada identifikasi metadata dokumen (seperti tanggal berlaku dan nomor surat) untuk memfilter dokumen relevan, bukan melakukan analisis interpretasi hukum mendalam terhadap konflik antar pasal.

        \item \textbf{Bahasa}: Sistem hanya dirancang untuk memproses dan menghasilkan jawaban dalam Bahasa Indonesia.

        \item \textbf{Input Pengguna}: Pertanyaan yang diajukan pengguna diasumsikan berupa pertanyaan terkait kebijakan universitas, bukan pertanyaan yang bersifat opini atau prediksi masa depan.
    \end{enumerate}

\section{Tujuan Penelitian}
    Tujuan utama dari penelitian ini adalah mengembangkan sistem pencarian informasi berbasis \textit{Retrieval-Augmented Generation} (RAG) yang mampu menangani perubahan regulasi universitas. 
    Secara spesifik, tujuan penelitian ini adalah:
    \begin{enumerate}
        \item Mengimplementasikan strategi \textit{chunking} dan \textit{embedding} yang mampu mempreservasi struktur hirarki dokumen (bab, pasal, ayat) untuk menjaga konteks semantik yang utuh.

        \item Merancang mekanisme sitasi otomatis dalam proses \textit{generation} untuk memastikan setiap jawaban yang diberikan memiliki rujukan ke dokumen regulasi yang digunakan.

        \item Mengevaluasi performa sistem RAG yang dikembangkan menggunakan metrik relevansi dan akurasi (seperti \textit{Precision}, \textit{Recall}, atau \textit{RAGAS score}) dibandingkan dengan pendekatan pencarian konvensional.
    \end{enumerate}

\section{Manfaat Penelitian}
Penelitian ini diharapkan dapat memberikan manfaat sebagai berikut:
\begin{enumerate}
    \item Membantu Sivitas Akademika Universitas Pertamina dalam mencari informasi regulasi secara akurat, khususnya dalam memastikan dokumen yang ditemukan adalah versi yang berlaku dan valid.

    \item Meningkatkan efisiensi pencarian dokumen hukum universitas melalui teknologi \textit{chatbot} berbasis \textit{Retrieval-Augmented Generation} (RAG) yang mampu memahami konteks pertanyaan pengguna dengan lebih baik dibandingkan pencarian kata kunci biasa.

    \item Menambah wawasan dan menjadi referensi bagi peneliti lain dalam pengembangan sistem tanya jawab otomatis (\textit{chatbot}) berbasis teknologi \textit{Retrieval-Augmented Generation} (RAG).
\end{enumerate}